\chapter{Développement de la solution adoptée}

\section{Introduction}

Après avoir présenté la conception globale de la solution, ce chapitre décrit le développement des principales composantes retenues dans le cadre du projet. Il ne s’agit pas uniquement de présenter un pipeline de traitement, mais d’expliquer la mise en place progressive d’une solution complète permettant de structurer les données, de les centraliser, d’appliquer certaines règles de validation et de les rendre accessibles à travers des interfaces de restitution.

Le développement de la solution repose d’abord sur un pipeline Data Engineering permettant de transformer les fichiers opérationnels \textit{Data Dump} et \textit{Coverage} en données structurées et exploitables. Ces deux fichiers constituent les principales sources utilisées dans le projet : le premier apporte les réponses détaillées issues des visites terrain, tandis que le second fournit une vision opérationnelle des visites planifiées, réalisées, rejetées, déviées ou non effectuées.

À partir de ces données, la solution construit progressivement une base MySQL organisée autour de tables de référence, de tables opérationnelles, de tables dynamiques par tâche et d’une table normalisée \textit{survey\_responses}. Cette structuration permet ensuite d’alimenter les analyses, les tableaux de bord, les règles de validation et les interfaces applicatives.

Ce chapitre présente également la mise en place des règles de validation métier, notamment la distinction entre les contrôles appliqués aux formats MT et GT. Enfin, il introduit la couche applicative développée autour de services backend Spring Boot, d’une interface mobile de consultation client et d’un portail back-office interne.

\section{Mise en place du pipeline Data Engineering}

La figure suivante présente une vue simplifiée du flux général du pipeline Data Engineering mis en place dans le cadre du projet. Ce flux met en évidence les principales étapes du traitement, depuis l’acquisition des fichiers Excel jusqu’au chargement des données structurées dans MySQL. Il montre également le rôle de Prefect dans l’orchestration du pipeline, ainsi que la logique progressive de transformation des données à travers les couches Bronze, Silver et Gold.

\begin{figure}[H]
    \centering
    \includegraphics[width=1\textwidth]{figures/etl/ETL.png}
    \caption{Flux général du pipeline Data Engineering}
    \label{fig:flux_etl}
\end{figure}

Le pipeline commence par les fichiers sources \textit{Data Dump} et \textit{Coverage}, qui peuvent être fournis à partir d’un chemin local ou récupérés via l’export du portail. Les données sont ensuite lues et préparées à l’aide de scripts Python, puis structurées progressivement. La couche Bronze correspond à l’ingestion initiale des fichiers, la couche Silver regroupe les opérations de nettoyage et d’harmonisation, tandis que la couche Gold prépare les tables prêtes à l’exploitation. Les résultats sont enfin chargés dans la base MySQL afin d’alimenter les règles de validation, les jeux de données BI et les services backend.

\subsection{Ingestion des fichiers Excel}

L’ingestion constitue la première étape du pipeline Data Engineering. Elle consiste à récupérer les fichiers opérationnels nécessaires au traitement, principalement le fichier \textit{Data Dump} et le fichier \textit{Coverage}. Ces deux fichiers représentent les sources de base du projet : le \textit{Data Dump} contient les réponses détaillées issues des visites terrain, tandis que le \textit{Coverage} apporte une vision opérationnelle des visites planifiées, réalisées, rejetées, déviées ou non effectuées.

Afin de rendre le traitement flexible, deux modes d’entrée ont été prévus. Le premier mode repose sur l’utilisation de fichiers Excel déjà disponibles localement. Dans ce cas, l’utilisateur peut lancer le traitement en indiquant directement le chemin du fichier \textit{Data Dump} et du fichier \textit{Coverage}. Ce mode est utile lorsque les exports ont déjà été téléchargés depuis le portail et doivent simplement être retraités par le pipeline.

Le second mode repose sur l’automatisation de l’export depuis le portail Smollan. Cette approche utilise Playwright afin d’automatiser certaines actions réalisées habituellement dans un navigateur, telles que l’accès au portail, l’application des paramètres nécessaires et le téléchargement des fichiers Excel. Ce mode permet de réduire les manipulations manuelles répétitives, tout en conservant les fichiers exportés comme point d’entrée du pipeline.

\subsubsection{Lancement du pipeline}

Le pipeline est lancé depuis le terminal à partir de l’environnement virtuel du projet. Cette étape permet d’exécuter le module principal \textit{main.py}, situé dans le dossier \textit{data-engineering}. Le lancement de cette commande ouvre ensuite l’interface de sélection, à partir de laquelle l’utilisateur peut choisir les fichiers à traiter et le mode d’entrée souhaité.

\begin{figure}[H]
    \centering
    \includegraphics[width=0.82\textwidth]{figures/etl/commande_lancement_pipeline.png}
    \caption{Commande de lancement du pipeline ETL}
    \label{fig:commande_lancement_pipeline}
\end{figure}

La figure \ref{fig:commande_lancement_pipeline} présente la commande utilisée pour démarrer le pipeline depuis le répertoire du projet. Cette commande constitue le point d’entrée de l’exécution avant la sélection des fichiers \textit{Data Dump} et \textit{Coverage}.

\subsubsection{Interface de sélection des fichiers}

Après le lancement de la commande, une interface graphique permet de configurer l’exécution du pipeline. Cette interface distingue les fichiers journaliers à traiter, à savoir le \textit{Data Dump} et le \textit{Coverage}. Elle permet également de choisir la source des données : extraction depuis le portail Smollan ou chargement de fichiers Excel présents localement sur la machine.

\begin{figure}[H]
    \centering
    \includegraphics[width=0.95\textwidth]{figures/etl/etl_launcher_interface.png}
    \caption{Interface de sélection des fichiers du pipeline ETL}
    \label{fig:etl_launcher_interface}
\end{figure}

La figure \ref{fig:etl_launcher_interface} montre les principales options disponibles au lancement. Le mode portail permet d’extraire les fichiers sélectionnés avant leur chargement dans MySQL, tandis que le mode local permet de sélectionner directement les fichiers Excel depuis la machine. L’interface propose également l’exécution des règles de validation après le chargement du \textit{Data Dump}, ce qui permet d’intégrer les contrôles métier dans le flux de traitement.

\subsubsection{Configuration de l’environnement}

La configuration du pipeline repose sur un fichier \textit{.env}. Ce fichier permet de centraliser les paramètres nécessaires à la connexion à la base MySQL, à l’export depuis le portail Smollan et à l’organisation des répertoires utilisés par le pipeline. Les informations sensibles, notamment les identifiants et les mots de passe, ne sont pas affichées dans le rapport et doivent rester anonymisées.

\begin{figure}[H]
    \centering
    \includegraphics[width=0.82\textwidth]{figures/etl/config_env_pipeline.png}
    \caption{Exemple anonymisé de configuration de l’environnement}
    \label{fig:config_env_pipeline}
\end{figure}

La figure \ref{fig:config_env_pipeline} montre les principaux paramètres utilisés par le pipeline. Les variables liées à MySQL permettent d’identifier la base \textit{unilever\_db}. Les variables liées au portail définissent l’URL, l’entité concernée, les rapports à extraire et les options d’exécution de Playwright. Enfin, les chemins de téléchargement, de sauvegarde et de logs permettent d’organiser les fichiers générés ou utilisés pendant le traitement.

\subsubsection{Orchestration et suivi d’exécution avec Prefect}

L’exécution du pipeline est orchestrée à l’aide de Prefect afin de mieux structurer les différentes étapes du traitement et de suivre leur état d’exécution. Dans le cadre du projet, Prefect ne sert pas uniquement à lancer le pipeline, mais à organiser son déroulement autour de phases fonctionnelles clairement identifiables : les contrôles préalables, l’acquisition des fichiers, le traitement du \textit{Data Dump}, l’exécution des règles de validation, le traitement du \textit{Coverage} et la synthèse finale de l’exécution.

Avant le lancement d’un traitement supervisé, le serveur Prefect doit être démarré localement. Cette étape permet d’activer l’API locale utilisée par le pipeline et de rendre disponible l’interface de suivi des exécutions. Lorsque le serveur n’est pas actif, le pipeline ne peut pas communiquer avec l’API de Prefect, ce qui empêche le suivi correct du run.

\begin{figure}[H]
\centering
\includegraphics[width=0.75\textwidth]{figures/etl/commande_prefect_server.png}
\caption{Démarrage du serveur Prefect}
\label{fig:commande_prefect_server}
\end{figure}

La figure \ref{fig:commande_prefect_server} présente la commande utilisée pour démarrer le serveur Prefect. Une fois le serveur lancé, Prefect met à disposition une API locale ainsi qu’un tableau de bord permettant de suivre les exécutions du pipeline, de consulter les runs et d’identifier l’état des différentes étapes.

La structure retenue repose sur un flux principal intitulé \textit{Merch Performance Daily ETL Orchestration}. Ce flux principal est organisé en sous-flux fonctionnels afin de séparer les responsabilités du pipeline sans afficher toutes les opérations internes comme des tâches indépendantes. Les opérations de bas niveau, telles que les transformations de \textit{DataFrames}, les insertions SQL, les boucles d’idempotence ou les requêtes propres aux règles de validation, restent encapsulées dans des fonctions Python internes.

\begin{figure}[H]
\centering
\includegraphics[width=0.95\textwidth]{figures/etl/prefect_runs_overview.png}
\caption{Vue des runs supervisés dans Prefect}
\label{fig:prefect_runs_overview}
\end{figure}

La figure \ref{fig:prefect_runs_overview} montre les principaux runs enregistrés dans Prefect. On y distingue le flux principal ainsi que les sous-flux associés aux grandes phases du traitement : contrôles préalables, acquisition et traçabilité des fichiers, traitement du \textit{Data Dump}, traitement des validations et traitement du \textit{Coverage}. Cette organisation rend le suivi plus lisible et permet d’identifier rapidement la phase concernée en cas d’erreur.

\begin{figure}[H]
\centering
\includegraphics[width=0.95\textwidth]{figures/etl/prefect_parent_flow.png}
\caption{Vue globale de l’orchestration du pipeline dans Prefect}
\label{fig:prefect_parent_flow}
\end{figure}

La figure \ref{fig:prefect_parent_flow} illustre la relation entre le flux principal et les sous-flux fonctionnels. Le pipeline commence par les vérifications préalables et l’acquisition des fichiers, puis enchaîne avec le traitement du \textit{Data Dump}, l’exécution des règles de validation, le traitement du \textit{Coverage} et la synthèse finale. Cette structuration permet de conserver un équilibre entre lisibilité dans l’interface Prefect et niveau de détail nécessaire au suivi technique du pipeline.

Ainsi, Prefect apporte une meilleure traçabilité par rapport à une exécution classique d’un script unique. Il permet de suivre les temps d’exécution, de distinguer les étapes fonctionnelles du pipeline et de faciliter l’analyse des erreurs éventuelles.

\subsection{Traitement du fichier Data Dump}

Le fichier \textit{Data Dump} constitue la source principale des réponses terrain. Il contient les informations issues des visites réalisées dans l’application Smart 3, notamment les données liées aux merchandisers, aux magasins, aux produits, aux tâches, aux questions et aux réponses saisies. Son traitement consiste à transformer ces lignes brutes en tables structurées dans MySQL.

Afin d’illustrer cette transformation, deux exemples représentatifs sont retenus : les réponses OSA et les données SOS. Ces deux cas montrent que le pipeline n’applique pas une transformation unique à toutes les tâches, mais adapte la structure SQL selon la logique métier de chaque type de donnée.

\begin{figure}[H]
    \centering
    \includegraphics[width=0.90\textwidth]{figures/etl/data_dump_to_sql_transformation.png}
    \caption{Transformation du Data Dump vers les tables SQL : exemples OSA et SOS}
    \label{fig:data_dump_to_sql_transformation}
\end{figure}

\subsubsection{Principe général de transformation}

Le traitement commence par la lecture des lignes du fichier \textit{Data Dump}. Chaque ligne correspond à une réponse terrain associée à une tâche, un titre, une question et une réponse. Le pipeline identifie ensuite la tâche métier concernée, rattache la ligne à la visite correspondante et, lorsque cela est possible, l’associe également au produit concerné.

Cette étape permet de passer d’une logique de fichier Excel, orientée lignes de réponses, vers une logique de tables SQL structurées. Les champs de liaison utilisés pour rattacher les réponses aux visites reposent principalement sur la date, le code merchandiser et le code magasin. Le code produit permet ensuite d’enrichir certaines tâches, notamment les tâches liées à la disponibilité produit.

\subsubsection{Étude de cas OSA : transformation en colonnes dynamiques}

Le premier cas concerne les réponses OSA. Dans le fichier \textit{Data Dump}, ces réponses sont stockées sous forme de lignes question/réponse associées à une visite et à un produit. Le traitement appliqué vise à transformer ces lignes brutes en une structure SQL plus adaptée à l’analyse de la disponibilité produit.

\begin{table}[H]
\centering

\begin{tcolorbox}[technicalbox,breakable=false,title={Étude de cas OSA --- Pivot dynamique},halign title=center]
\small
\renewcommand{\arraystretch}{1.35}
\begin{tabular}{p{0.23\textwidth}|p{0.67\textwidth}}

\textbf{Entrée Data Dump} &
Ligne question/réponse liée à une visite, à un produit et à une réponse de disponibilité du SKU, généralement sous la forme \textit{Oui} ou \textit{Non}. \\
\hline

\textbf{Règle appliquée} &
Les lignes dont le titre contient \textit{OSA}, \textit{COC}, \textit{MH} ou \textit{PACK} sont orientées vers la table \texttt{task\_osa\_pack\_coc\_mh}. \\
\hline

\textbf{Transformation} &
Les réponses sont regroupées par \texttt{visit\_id} et \texttt{product\_id}. Chaque question devient une colonne SQL dynamique, tandis que la réponse devient la valeur stockée dans cette colonne. \\
\hline

\textbf{Sortie SQL} &
La table cible contient une ligne par visite et par produit, avec des colonnes fixes comme \texttt{visit\_id}, \texttt{product\_id}, \texttt{title} et \texttt{response\_date}, ainsi que des colonnes dynamiques générées à partir des questions. \\
\hline

\textbf{Utilité} &
Cette structure facilite l’analyse de la disponibilité des produits et prépare les données pour les indicateurs OSA et les règles de contrôle métier.

\end{tabular}
\end{tcolorbox}

\caption{Étude de cas OSA : transformation en colonnes dynamiques}
\label{tab:osa_transformation}
\end{table}


Cette logique correspond à une transformation de type \textit{pivot}. Elle permet de passer de plusieurs lignes question/réponse vers une ligne structurée par produit visité.

\subsubsection{Étude de cas SOS : transformation en structure verticale}

Le second cas concerne les données SOS. Contrairement à l’OSA, les réponses SOS ne sont pas transformées en colonnes dynamiques. Elles sont conservées sous forme verticale afin de représenter les mesures de part de linéaire, les valeurs relevées et le total associé.

\begin{table}[H]
\centering

\begin{tcolorbox}[technicalbox,breakable=false,title={Étude de cas SOS --- Structure verticale},halign title=center]
\small
\renewcommand{\arraystretch}{1.35}
\begin{tabular}{p{0.23\textwidth}|p{0.67\textwidth}}

\textbf{Entrée Data Dump} &
Ligne liée à une mesure de part de linéaire, contenant une question, une valeur relevée et un contexte de visite. \\
\hline

\textbf{Règle appliquée} &
Les lignes dont le titre contient \textit{SOS} sont orientées vers la table \texttt{task\_sos}. Le traitement ne transforme pas les questions en colonnes dynamiques. \\
\hline

\textbf{Transformation} &
La question est conservée dans \texttt{qst}, la réponse dans \texttt{qst\_value}, et le total associé dans \texttt{total}. Chaque ligne SOS reste donc un item d’analyse. \\
\hline

\textbf{Sortie SQL} &
La table cible contient une ligne par item SOS, avec les champs \texttt{visit\_id}, \texttt{task}, \texttt{response\_date}, \texttt{qst}, \texttt{qst\_value} et \texttt{total}. \\
\hline

\textbf{Utilité} &
Cette structure facilite les agrégations et les comparaisons nécessaires à l’analyse de la part de linéaire et de la visibilité en rayon.

\end{tabular}
\end{tcolorbox}

\caption{Étude de cas SOS : transformation en structure verticale}
\label{tab:sos_transformation}
\end{table}

Cette logique verticale est plus adaptée au SOS, car elle conserve les mesures sous une forme directement exploitable pour les analyses quantitatives.

\subsubsection{Sorties générées et fiabilisation du rechargement}

À partir du fichier \textit{Data Dump}, le pipeline alimente plusieurs familles de tables. Les tables \texttt{employees}, \texttt{stores}, \texttt{products} et \texttt{visits} structurent les principales entités métier. Les tables \texttt{task\_*} conservent les réponses selon leur logique opérationnelle, tandis que \texttt{survey\_responses} fournit une vue normalisée des réponses terrain.

Afin de fiabiliser les relances du pipeline, les données détaillées associées aux visites retraitées sont nettoyées avant leur réinsertion. Cette opération concerne les tables \texttt{task\_*} et \texttt{survey\_responses}. Elle permet d’éviter l’accumulation de doublons et de garantir que les données détaillées correspondent au dernier fichier traité.

\subsection{Intégration du Coverage et rapprochement avec le Data Dump}

Le fichier \textit{Coverage} constitue la deuxième source principale du pipeline. Contrairement au fichier \textit{Data Dump}, qui contient les réponses détaillées issues des tâches terrain, le \textit{Coverage} apporte une vision opérationnelle de l’exécution des visites. Il permet notamment de suivre les visites planifiées, réalisées, rejetées, déviées ou non effectuées.

Dans la solution, le \textit{Coverage} est chargé dans la table \texttt{fact\_coverage}. Cette table conserve le grain opérationnel du fichier source, basé principalement sur la date de visite, le merchandiser et le magasin. Cette séparation est volontaire, car certaines lignes du \textit{Coverage}, comme les visites non effectuées ou rejetées, peuvent exister sans disposer de réponses détaillées dans le \textit{Data Dump}.

\subsubsection{Rôle de la table de couverture opérationnelle}

La table \texttt{fact\_coverage} joue le rôle d’une table de faits opérationnelle. Elle centralise les informations nécessaires au suivi de l’exécution terrain : statut de visite, couverture, rejet, déviation, non-visite, distances GPS, temps d’exécution et motifs éventuels.

Son objectif n’est pas de remplacer les tables issues du \textit{Data Dump}, mais de compléter leur lecture. Le \textit{Data Dump} permet d’analyser le contenu détaillé des visites exécutées, tandis que le \textit{Coverage} permet d’évaluer l’état global de l’exécution, y compris les cas où aucune réponse terrain détaillée n’est disponible.

\subsubsection{Rapprochement optionnel avec les visites détaillées}

Afin de relier le suivi opérationnel issu du \textit{Coverage} avec les visites détaillées issues du \textit{Data Dump}, une vue SQL de rapprochement a été ajoutée : \texttt{vw\_fact\_coverage\_visit\_match}. Cette vue ne remplace pas les tables existantes et ne modifie pas leur structure. Elle agit comme une couche logique de lecture permettant de comparer les lignes de \texttt{fact\_coverage} avec les visites présentes dans la table \texttt{visits}.

L’intérêt de cette approche est de conserver toutes les lignes du \textit{Coverage}, y compris les visites non effectuées, rejetées ou planifiées sans données détaillées. Le rapprochement est donc volontairement optionnel : lorsqu’une visite détaillée existe dans le \textit{Data Dump}, la vue expose son \texttt{visit\_id}. Dans le cas contraire, la ligne Coverage reste visible avec un statut indiquant l’absence de visite détaillée.

\begin{figure}[H]
    \centering
    \includegraphics[width=0.37\textwidth]{figures/database/view_coverage_datadump_object.png}
    \caption{Implémentation de la vue SQL de rapprochement dans la base \texttt{unilever\_db}}
    \label{fig:view_coverage_datadump_object}
\end{figure}

La figure \ref{fig:view_coverage_datadump_object} montre que le rapprochement a été implémenté sous forme d’une vue SQL distincte, et non comme une nouvelle table physique. Cette séparation permet de garder \texttt{fact\_coverage} comme table de faits opérationnelle, tout en ajoutant une lecture croisée avec les visites détaillées lorsque celles-ci existent.

\begin{tcolorbox}[technicalbox,title={Logique de la vue de rapprochement}]
\small

\textbf{Source principale :} \texttt{fact\_coverage}, qui conserve les statuts opérationnels des visites : visité, non visité, rejeté ou en déviation.

\vspace{0.12cm}

\textbf{Source rapprochée :} \texttt{visits}, alimentée à partir du \textit{Data Dump} et représentant les visites disposant de données détaillées.

\vspace{0.12cm}

\textbf{Clé utilisée :} \texttt{visit\_date} + \texttt{employee\_code} + \texttt{store\_code}.

\vspace{0.12cm}

\textbf{Résultat ajouté par la vue :} le champ \texttt{visit\_id}, lorsqu’une correspondance est trouvée, et le champ \texttt{match\_status} pour qualifier le rapprochement.

\vspace{0.12cm}

\textbf{Statut \texttt{MATCHED} :} la ligne Coverage possède une visite détaillée correspondante dans le \textit{Data Dump}.

\vspace{0.12cm}

\textbf{Statut \texttt{NO\_DATA\_DUMP\_VISIT} :} la ligne existe dans le \textit{Coverage}, mais aucune visite détaillée correspondante n’existe dans le \textit{Data Dump}.
\end{tcolorbox}

Cette approche renforce l’exploitation du fichier \textit{Coverage}. Celui-ci n’est pas seulement chargé dans une table, mais intégré comme une source de suivi opérationnel pouvant être comparée aux visites détaillées issues du \textit{Data Dump}. La vue permet ainsi d’identifier les visites disposant de réponses terrain, tout en conservant les situations présentes uniquement dans le suivi opérationnel.

\subsubsection{Apport dans la solution}

Cette modélisation renforce la cohérence globale de la solution. Elle évite de forcer une relation physique entre \texttt{fact\_coverage} et \texttt{visits}, ce qui aurait pu poser problème pour les visites non effectuées ou rejetées. À la place, la vue SQL propose une lecture contrôlée et réversible, sans imposer de modification directe aux tables existantes ni aux services applicatifs déjà mis en place.

Elle constitue également une base utile pour les analyses futures, car elle distingue clairement les visites ayant des réponses détaillées de celles qui existent uniquement dans le suivi de couverture. Cette distinction peut ensuite être exploitée dans les tableaux de bord, les contrôles opérationnels ou les interfaces de suivi.

\subsection{Modélisation et chargement dans MySQL}

Après les étapes de préparation et de transformation, les données sont chargées dans une base MySQL nommée \textit{unilever\_db}. Cette base constitue le référentiel central de la solution. Elle permet de conserver les données issues des fichiers \textit{Data Dump} et \textit{Coverage}, de structurer les informations métier et de les rendre accessibles aux règles de validation, aux tableaux de bord et aux services applicatifs.

Le modèle retenu ne cherche pas à reproduire directement la structure des fichiers Excel. Il organise les données selon leur rôle dans la solution : les entités de référence, les visites, les réponses détaillées, les données de couverture, les résultats de validation et les journaux d’exécution.

\begin{figure}[H]
    \centering
    \includegraphics[width=1\textwidth]{figures/database/mysql_logical_model.png}
    \caption{Organisation logique des données dans MySQL}
    \label{fig:mysql_logical_model}
\end{figure}

La figure \ref{fig:mysql_logical_model} présente une vue simplifiée de l’organisation des données dans MySQL. Les fichiers \textit{Data Dump} et \textit{Coverage} alimentent des tables ayant des rôles complémentaires. Le \textit{Data Dump} permet de construire les tables de référence, les visites, les tables de tâches et la table normalisée \texttt{survey\_responses}. Le \textit{Coverage} alimente la table \texttt{fact\_coverage}, utilisée pour le suivi opérationnel des visites.

Cette organisation permet de séparer les données stables, les données opérationnelles et les données analytiques. Elle facilite également l’exploitation des données par les règles de validation, les tableaux de bord et les APIs applicatives.

\subsubsection{Stratégie de chargement}

Le chargement des données s’effectue de manière contrôlée afin d’éviter les doublons et de garantir la cohérence des informations insérées. Les tables de référence sont alimentées à partir des codes métier, comme le code employé, le code magasin ou le code produit. Lorsque ces éléments existent déjà dans la base, les informations peuvent être mises à jour sans créer de doublons.

Les visites sont chargées à partir d’une clé naturelle composée de la date de visite, du merchandiser et du magasin. Cette clé permet d’identifier une visite de manière cohérente et de relier ensuite les réponses terrain aux visites correspondantes.

Pour les données détaillées issues du \textit{Data Dump}, le pipeline applique une logique de rechargement contrôlé. Avant d’insérer les nouvelles réponses associées aux visites retraitées, les anciennes lignes correspondantes sont supprimées des tables de tâches et de la table \texttt{survey\_responses}. Cette étape permet de reconstruire proprement les données détaillées à partir du dernier fichier traité.

\subsubsection{Traçabilité des traitements}

La traçabilité constitue un élément important du chargement. Le pipeline conserve des informations sur les exécutions réalisées, notamment à travers la table \texttt{etl\_run\_log}. Cette table permet de suivre les traitements effectués, les fichiers concernés, les dates d’exécution et l’état global du chargement.

Cette logique de suivi facilite l’analyse en cas d’erreur ou d’incohérence. Elle permet également de mieux comprendre l’historique des traitements appliqués aux données et de vérifier que les différentes étapes du pipeline ont été exécutées correctement.

Ainsi, le chargement dans MySQL représente une étape centrale de la solution. Il transforme les fichiers Excel traités en un modèle relationnel exploitable, capable d’alimenter les validations métier, les tableaux de bord et les interfaces de consultation.

\section{Structuration de la base de données}

La base MySQL \textit{unilever\_db} est structurée selon le rôle des données dans la solution. L’objectif n’est pas de reproduire directement les fichiers Excel, mais d’organiser les informations autour de plusieurs familles de tables : tables de référence, tables opérationnelles, tables détaillées d’exécution, vue de rapprochement, tables de validation et tables de traçabilité.

\begin{figure}[H]
\centering
\includegraphics[width=1\textwidth]{figures/database/tables\_flow.png}
\caption{Modèle relationnel simplifié de la base \texttt{unilever\_db}}
\label{fig:dbdiagram_simplified_model}
\end{figure}

La figure \ref{fig:dbdiagram_simplified_model} présente un modèle relationnel simplifié de la base. Elle ne reprend pas toutes les colonnes techniques, mais met en évidence les principales relations entre les visites, les magasins, les produits, les réponses terrain, la couverture opérationnelle et les résultats de validation.

Les tables \texttt{employees}, \texttt{stores} et \texttt{products} jouent le rôle de tables de référence. Elles décrivent les entités métier principales et évitent de répéter les mêmes informations dans les tables d’exécution. Les tables \texttt{visits} et \texttt{fact\_coverage} représentent les données opérationnelles : la première porte les visites détaillées issues du \textit{Data Dump}, tandis que la seconde conserve le suivi opérationnel issu du \textit{Coverage}.

Les tables \texttt{task\_*} stockent les réponses détaillées selon la nature des tâches terrain. Certaines tâches, comme l’OSA, utilisent une logique de colonnes dynamiques, tandis que d’autres, comme le SOS, conservent une structure verticale. La table \texttt{survey\_responses} complète ces tables spécialisées en fournissant une vue normalisée et transversale des réponses terrain.

La vue \texttt{vw\_fact\_coverage\_visit\_match} permet de rapprocher le \textit{Coverage} et les visites détaillées sans imposer de relation physique obligatoire. Enfin, les tables \texttt{validation\_rules}, \texttt{validation\_run\_log}, \texttt{validation\_results} et \texttt{etl\_run\_log} assurent le suivi des règles métier, des résultats de validation et des exécutions du pipeline.

\section{Mise en place du moteur de validation}

Après la structuration des données dans MySQL, une couche de validation a été mise en place afin de détecter certaines situations nécessitant une vérification opérationnelle. Ce moteur est implémenté en Python et s’appuie sur les données déjà chargées dans la base, notamment les visites, les magasins, les merchandisers et les réponses terrain.

L’objectif de cette couche n’est pas de remplacer le contrôle humain, mais d’aider à identifier des situations inhabituelles. Les règles appliquées permettent de repérer des cas pouvant nécessiter une revue, comme des réponses OSA incohérentes avec le comportement global d’une bannière ou des coordonnées GPS inhabituelles sur des visites répétées.

\begin{tcolorbox}[technicalbox,title={Synthèse des règles implémentées}]
\small
\textbf{Règle OSA -- formats MT :}
détection d’une réponse \textit{Non} inhabituelle lorsque le comportement hebdomadaire de la même bannière et du même produit indique une disponibilité majoritairement positive. La règle s’appuie principalement sur \texttt{survey\_responses}.

\vspace{0.18cm}
\hrule
\vspace{0.18cm}
\textbf{Règle GPS -- formats GT :}
détection d’une visite dont la position GPS s’éloigne fortement de la position médiane observée pour le même merchandiser et le même magasin durant le mois. La règle s’appuie principalement sur \texttt{visits}.
\end{tcolorbox}

\subsection{Principe général du moteur de validation}

Le moteur de validation est organisé autour d’un registre de règles Python. Les règles disponibles sont déclarées dans ce registre, puis synchronisées avec la table \texttt{validation\_rules}. Cette table permet de conserver les règles connues par le système et d’indiquer si elles sont actives ou non à l’aide du champ \texttt{is\_active}.

Lorsqu’une exécution de validation démarre, une ligne est créée dans \texttt{validation\_run\_log}. Le moteur charge ensuite les règles actives, exécute les contrôles correspondants et insère les situations détectées dans \texttt{validation\_results}. À la fin du traitement, le journal d’exécution est mis à jour avec le statut final, le nombre de règles exécutées et le nombre de situations détectées.

Dans l’orchestration actuelle, la validation est exécutée après le traitement du fichier \textit{Data Dump}. Elle s’appuie donc principalement sur les tables alimentées par ce fichier, notamment \texttt{visits} et \texttt{survey\_responses}. L’exécution peut aussi être limitée aux visites affectées par le chargement en cours, afin d’éviter de retraiter inutilement l’ensemble de l’historique.

\subsection{Règle de validation OSA pour les formats MT}

La première règle concerne les réponses OSA pour les formats MT. Elle s’applique uniquement aux magasins de type \textit{HYPERMARKET}, \textit{SUPERMARKET LOWER} et \textit{SUPERMARKET UPPER}. Son objectif est de repérer des réponses de disponibilité produit qui paraissent inhabituelles par rapport au comportement observé sur une même bannière.

La règle analyse les réponses \textit{Oui} et \textit{Non} dans \texttt{survey\_responses}, puis les rapproche avec les visites et les magasins afin de prendre en compte le format du point de vente et la bannière. Une situation est signalée lorsqu’une réponse individuelle indique \textit{Non}, alors que le comportement hebdomadaire de la même bannière, du même produit et de la même question montre un taux de disponibilité élevé.

Dans l’implémentation actuelle, une situation est générée lorsque le taux de disponibilité hebdomadaire est supérieur ou égal à 90\%.

\subsection{Règle de validation GPS pour les formats GT}

La seconde règle concerne la cohérence GPS des visites répétées dans les formats GT. Elle s’applique uniquement aux magasins de type \textit{GROCERY} et \textit{CASH AND CARRY}. Son objectif est d’identifier des coordonnées GPS inhabituelles pour des visites répétées au même point de vente.

La règle analyse les coordonnées enregistrées dans \texttt{visits}. Les visites sont regroupées par merchandiser, magasin, année et mois. Pour chaque groupe disposant d’un nombre suffisant de visites, le moteur calcule une position de référence à partir des coordonnées médianes, puis compare chaque visite à cette position.

Une situation est signalée lorsque la distance entre la position GPS d’une visite et la position médiane du groupe dépasse un seuil défini. Dans l’implémentation actuelle, une anomalie est créée lorsque cette distance dépasse 1000 mètres. La sévérité devient plus élevée lorsque la distance atteint un seuil plus important.

\subsection{Stockage des résultats de validation}

Les résultats produits par les règles sont stockés dans la table \texttt{validation\_results}. Cette table conserve les informations nécessaires à l’analyse des situations détectées : règle appliquée, visite concernée, magasin, merchandiser, produit lorsque celui-ci est disponible, valeur observée, valeur attendue, métrique calculée, message de validation et niveau de sévérité.

La table contient également des champs de revue, tels que le statut de traitement, le commentaire, l’utilisateur ayant réalisé la revue et la date de revue. Ces champs permettent d’exploiter les situations détectées dans un portail back-office, où elles peuvent être consultées, analysées et suivies.

La table \texttt{validation\_run\_log} assure la traçabilité des exécutions du moteur. Elle permet de savoir quand une validation a été lancée, si elle s’est terminée correctement, combien de règles ont été exécutées et combien de situations ont été détectées.

Dans l’état actuel de la solution, le moteur de validation reste limité à deux règles principales. Les seuils et les formats concernés sont définis dans le code Python. Cette approche reste adaptée au périmètre actuel du projet, tout en laissant la possibilité d’évoluer vers un moteur plus configurable.

\section{Développement de la couche Business Intelligence}

La couche Business Intelligence constitue la partie analytique de la solution. Elle permet de transformer les données centralisées dans MySQL en indicateurs de suivi opérationnel exploitables par les équipes métier. Son rôle est de fournir une lecture synthétique de l’exécution terrain, en complément des interfaces applicatives et des règles de validation.

Dans le cadre du projet, Power BI a été utilisé comme outil de restitution afin de construire une première version analytique autour de la couverture, de la performance des merchandisers, des visites problématiques et de certains indicateurs liés à l’exécution terrain. Cette couche ne doit pas être considérée comme une plateforme BI totalement industrialisée, mais comme un modèle analytique préparé à partir de la base \texttt{unilever\_db}.

\subsection{Objectif de la couche BI}

L’objectif de la couche BI est de rendre les données opérationnelles plus lisibles pour le suivi métier. Les fichiers \textit{Data Dump} et \textit{Coverage} contiennent une grande quantité d’informations, mais leur exploitation directe reste limitée lorsqu’il s’agit de suivre plusieurs indicateurs, de filtrer par période, superviseur ou merchandiser, ou de comparer les performances d’exécution.

Power BI permet donc de regrouper les indicateurs clés dans des vues de pilotage. Ces vues facilitent le suivi du taux de couverture, de la complétion des tâches, des déviations, des visites non effectuées, des rejets, des distances GPS et de la performance par merchandiser. Elles permettent également de passer d’une lecture fichier par fichier à une lecture plus synthétique et interactive de l’exécution terrain.

\subsection{Connexion directe à la base MySQL}

La préparation de la couche BI commence par une connexion directe entre Power BI et la base MySQL \texttt{unilever\_db}. Ce choix permet d’exploiter les données déjà centralisées par le pipeline Data Engineering, au lieu de repartir des fichiers Excel bruts. La couche BI s’appuie donc sur le même référentiel de données que le reste de la solution.

\begin{figure}[H]
\centering
\begin{minipage}{0.42\textwidth}
\centering
\includegraphics[width=\textwidth]{figures/dashboards/powerbi_mysql_connector.png}
\end{minipage}
\hfill
\begin{minipage}{0.42\textwidth}
\centering
\includegraphics[width=\textwidth]{figures/dashboards/powerbi_mysql_connection.png}
\end{minipage}
\caption{Connexion de Power BI à la base MySQL \texttt{unilever\_db}}
\label{fig:powerbi_mysql_connection}
\end{figure}

La figure \ref{fig:powerbi_mysql_connection} illustre les premières étapes de connexion à la base MySQL. Power BI permet d’abord de sélectionner le connecteur MySQL, puis de renseigner les informations nécessaires à l’accès à la base locale, notamment le serveur, le port, le nom de la base et les identifiants de connexion. Cette connexion directe permet d’utiliser les données structurées dans MySQL comme source principale du modèle BI.

\subsection{Sélection des tables utiles}

Après la connexion à MySQL, les tables nécessaires sont sélectionnées depuis le navigateur Power BI. Toutes les tables de la base ne sont pas utilisées dans le modèle BI. La sélection se fait selon les besoins d’analyse : suivi de couverture, performance merchandiser, analyse des tâches terrain, filtres par magasin, période ou superviseur.

Les tables sélectionnées concernent principalement la table \texttt{fact\_coverage}, les visites, les référentiels métier et les tables de tâches. Les tables de validation, telles que \texttt{validation\_results}, \texttt{validation\_rules} et \texttt{validation\_run\_log}, sont plutôt destinées au portail back-office de suivi des situations détectées. Elles ne constituent donc pas la source principale de la couche BI opérationnelle.

\subsection{Préparation du modèle Power BI}

Une fois les tables importées, le modèle Power BI est structuré à travers les relations nécessaires entre les tables. Cette étape permet d’assurer la cohérence des filtres entre les dates, les merchandisers, les magasins et les données opérationnelles. Elle constitue une phase importante, car les indicateurs ne dépendent pas uniquement des tables importées, mais aussi de la manière dont elles sont reliées dans le modèle.

La table \texttt{fact\_coverage} constitue la source principale des indicateurs opérationnels, car elle contient les informations liées aux visites planifiées, réalisées, non visitées, rejetées ou en déviation. Elle est complétée par des tables de référence comme les employés, les magasins, les produits et les superviseurs. Des tables analytiques sont également préparées dans Power BI, comme \texttt{Date\_Dim}, \texttt{Dim\_Merch}, \texttt{Merch}, \texttt{Store}, \texttt{Store GPS Median}, \texttt{Executive} et \texttt{PREVIOUS PERIOD KPI}.

\begin{figure}[H]
\centering
\includegraphics[width=0.82\textwidth]{figures/dashboards/powerbi_model_relationships.png}
\caption{Modèle relationnel préparé dans Power BI}
\label{fig:powerbi_model_relationships}
\end{figure}

La figure \ref{fig:powerbi_model_relationships} présente le modèle relationnel préparé dans Power BI. Ce modèle met en relation les tables importées depuis MySQL avec des tables analytiques créées pour faciliter le filtrage temporel, l’analyse par merchandiser, l’analyse par magasin et le calcul des indicateurs. Cette étape permet de passer d’un simple import de tables à un modèle exploitable pour la restitution décisionnelle.

\subsection{Mesures et indicateurs construits}

Plusieurs mesures ont été construites dans Power BI afin de suivre l’exécution terrain. Les principaux indicateurs concernent le taux de couverture, la complétion des tâches, le taux de déviation, le taux de rejet, les visites non effectuées, les visites problématiques, la durée moyenne de visite et les indicateurs liés à la cohérence GPS.

Un indicateur composite de productivité a également été défini afin de résumer la performance d’un merchandiser à partir de plusieurs dimensions opérationnelles. Cet indicateur combine des facteurs positifs, tels que la couverture, la complétion des tâches et la conformité GPS, avec des facteurs pénalisants, comme les déviations et les rejets.

\begin{tcolorbox}[technicalbox,title={Logique du score de productivité}]
\small
Le score de productivité combine trois facteurs positifs : le taux de couverture, la complétion des tâches et la conformité GPS. Il applique également une pénalisation liée au taux de déviation et au taux de rejet. Cet indicateur a été construit dans le cadre du projet afin de faciliter la comparaison analytique entre merchandisers, sans remplacer l’analyse détaillée des indicateurs qui le composent.
\end{tcolorbox}

\subsection{Vues de restitution préparées}

Deux vues principales ont été préparées dans Power BI. La première correspond au \textit{National Retail Execution Dashboard}. Elle permet de suivre les indicateurs globaux de l’exécution terrain, notamment la couverture, les tâches réalisées, les visites planifiées, les visites non effectuées, les rejets, les déviations et les visites problématiques.

La seconde vue correspond au \textit{Merchandiser Performance Intelligence}. Elle est orientée vers l’analyse individuelle des merchandisers. Elle permet de suivre la couverture, la complétion des tâches, la durée de travail, la durée moyenne de visite, les indicateurs GPS, les visites non effectuées et les situations problématiques associées à un merchandiser donné.

Dans ce chapitre, la couche BI est présentée sous l’angle de sa préparation technique : connexion à MySQL, sélection des tables utiles, construction du modèle, création des mesures et préparation des vues de restitution. Les captures et l’interprétation des tableaux de bord seront présentées dans le chapitre consacré à la réalisation et aux résultats.

\section{Développement de la couche backend}

\subsection{Architecture Spring Boot et Java}
\subsection{APIs REST destinées à l’application mobile}
\subsection{APIs REST destinées au back-office}

\section{Développement des interfaces applicatives}

\subsection{Interface mobile Expo / React Native}
\subsection{Portail back-office React}
\subsection{Limites actuelles de l’authentification}

\section{Conclusion}
